% Lettre de motivation — Analyste Quantitatif, eXalt-Fi, Paris
% Compilation : xelatex lettre-exalt.tex
\documentclass[11pt,a4paper]{article}
\usepackage{fontspec}
\usepackage[french]{babel}
\usepackage[top=1.6cm, bottom=1.6cm, left=2.0cm, right=2.0cm]{geometry}
\usepackage{amssymb}
\usepackage{xcolor}
\usepackage{enumitem}
\usepackage{titlesec}
\usepackage[hidelinks]{hyperref}
\usepackage{fontawesome5}
\setmainfont{texgyrepagella}[Extension=.otf, UprightFont=*-regular, BoldFont=*-bold, ItalicFont=*-italic, BoldItalicFont=*-bolditalic]
\newfontfamily\headingfont{texgyreheros}[Extension=.otf, UprightFont=*-regular, BoldFont=*-bold, ItalicFont=*-italic, BoldItalicFont=*-bolditalic]
\definecolor{accent}{RGB}{28,58,94}
\definecolor{ink}{RGB}{26,26,26}
\definecolor{muted}{RGB}{95,105,120}
\definecolor{rulegray}{RGB}{200,206,214}
\color{ink}
\pagestyle{empty}
\setlength{\parindent}{0pt}
\setlength{\parskip}{0.6ex}
\linespread{1.03}
\titleformat{\section}{\headingfont\bfseries\fontsize{10}{12}\selectfont\color{accent}}{}{0em}{\MakeUppercase}[\vspace{-0.9ex}{\color{rulegray}\rule{\linewidth}{0.7pt}}]
\titlespacing*{\section}{0pt}{2.0ex plus .3ex}{1.0ex}
\titleformat{\paragraph}[runin]{\headingfont\bfseries\small\color{accent}}{}{0em}{}
\titlespacing*{\paragraph}{0pt}{1.0ex}{0.6em}
\setlist{itemsep=0.25ex, topsep=0.4ex, leftmargin=1.3em, label={\textcolor{accent}{\scriptsize$\blacktriangleright$}}}
\newcommand{\letterhead}{%
{\headingfont\bfseries\fontsize{16}{19}\selectfont EL MEHDI HARESS}\\[0.4ex]
{\footnotesize\color{muted}Inria Sophia Antipolis \textemdash{} CALISTO team\quad
\faEnvelope[regular]~\href{mailto:el-mehdi.haress@inria.fr}{el-mehdi.haress@inria.fr}\quad
\faPhone*~+33 7 58 23 41 50\quad
\faGlobe~\href{https://elmehdiharess.org}{elmehdiharess.org}}\\[-0.6ex]
{\color{rulegray}\rule{\linewidth}{0.7pt}}\par\vspace{1.2ex}}

\frenchspacing
\begin{document}
\letterhead

\begin{flushright}\small Sophia Antipolis, le \today\end{flushright}

\textbf{À l'attention de :} eXalt-Fi \textemdash{} équipe de recrutement Finance de Marché\\
\textbf{Objet :} Candidature au poste d'Analyste Quantitatif \textemdash{} Finance de Marché (CDI, Paris)

Madame, Monsieur,

Docteur en mathématiques appliquées de l'Université Paris-Saclay et ingénieur CentraleSupélec, je vous
adresse ma candidature au poste d'Analyste Quantitatif en Finance de Marché. Après un post-doctorat à
l'Université de Leeds, je suis aujourd'hui ingénieur de recherche à Inria Sophia Antipolis, et je souhaite
mettre mon expertise en modélisation stochastique au service des acteurs de la finance de marché.

Les missions que vous décrivez reposent sur les outils avec lesquels je travaille depuis cinq ans.
\begin{itemize}
\item \textbf{Modèles et pricing.} Équations différentielles stochastiques, schémas de discrétisation et
méthodes de Monte-Carlo avec garanties de convergence, y compris pour des dynamiques irrégulières où les
méthodes standard échouent (\emph{Stochastic Processes and their Applications}, \emph{IMA Journal of
Numerical Analysis}). J'ai également enseigné la finance stochastique et la modélisation du risque en M2
à CentraleSupélec.
\item \textbf{Risques et validation de modèles.} Mes travaux à Leeds quantifient la sensibilité d'une
dynamique aux perturbations de sa spécification et de ses paramètres, ce qui est au cœur de l'analyse du
risque de modèle. À l'Université d'Osaka, j'ai travaillé sur l'optimisation de mesures de risque
(VaR, CVaR, mesures spectrales) sous queues lourdes, avec deux articles à la conférence \emph{AISTATS}.
\item \textbf{Calibration et analyse statistique.} J'ai développé des estimateurs conjoints de plusieurs
paramètres (dérive, volatilité, mémoire) à partir d'observations discrètes, notamment pour un modèle de
taux d'Ornstein--Uhlenbeck fractionnaire, avec vitesses de convergence prouvées.
\item \textbf{Implémentation et documentation.} Je code en Python au quotidien depuis huit ans et je publie
mon code de recherche en \emph{open source}. La rédaction de notes techniques et la présentation de résultats
à des interlocuteurs non spécialistes font partie de mon travail, aujourd'hui auprès d'ingénieurs et de
partenaires industriels à Inria.
\end{itemize}

Je n'ai pas encore d'expérience en salle de marché, et c'est précisément ce qui m'attire dans le conseil :
intervenir auprès de banques, de gestionnaires d'actifs et de fintechs, sur des problématiques variées,
pour appliquer cette rigueur au pricing, à la mesure des risques de marché et de contrepartie et à la
validation de modèles. Mes bases en \mbox{C++} sont celles d'une formation d'ingénieur, et je suis prêt à
les renforcer rapidement selon les besoins de vos missions.

Mobile sur Paris, je serais heureux d'échanger avec vous sur la manière dont mon profil peut répondre aux
besoins de vos clients.

Je vous prie d'agréer, Madame, Monsieur, l'expression de mes salutations distinguées.

\vspace{0.6ex}
El Mehdi Haress

\end{document}
